\section{Business / Objective Understanding}
If a rider could close their eyes while cycling (which is not recommended by obvious reasons), this person's body would be able to clearly feel that they are passing either in a smooth or in a bumpy road. Perhaps one of the easiest ways to explain that in physics-related terms would be that when one is crossing a bumpy path, the dynamic vertical acceleration of one's bicycle is much more noticeable (and uncomfortable) if compared to a completely smooth terrain.

This definition gives us a strong hypothesis to test. If this very distinctive vibration feeling that the riders experiment can be derived from the sensors already present in a smartphone, for instance, that would be a clear and well-defined way to feed an autonomous classifier to ingest and analyse the sensor data, providing a real-time diagnostic whether a bike path is smooth or bumpy.

Naturally, this diagnostic is irrelevant on a user by user basis. As we already established, it's quite easy to feel the condition of the road you are cycling in. Accordingly, the value of such a classifier would be perceived for stakeholders responsible for providing adequate road conditions, or even more so, to equip the community with a tool to evaluate which parts of a given municipality need smoother cycling paths.

Considering the actual classifier, what is important to determine and establish in terms of scope is the ability of the model to separate between two distinct classes: \textit{smooth} and \textit{bumpy}. So, let's define that a single pothole in a smooth bicycle lane \textbf{should not} necessarily trigger a bumpy classification. As much as a pothole needs fixing to avoid accidents, this should fall in routine repairs and, as stated, the objective here is to raise awareness to bumpy lanes that need to be, at some point, converted to smooth ones.

\textbf{Research objective.}
As a high level model design criteria, the targets were defined as: $"Smooth" = 0$ and $"Bumpy" = 1$. Since the main objective is to locate paths which are bumpy across cities and villages, what is most important for this model is to avoid classifying smooth roads as bumpy, in other words, we need a model focused on delivering $\text{high precision}$ bumpy outcomes, supported by a low $\frac{\text{false positive rate}}{specificity}$, where a false positive is defined as: $\text{actual Smooth → predicted Bumpy}$.
\begin{itemize}
    \item headline metric      = Bumpy precision
    \item direct FP metric     = Smooth specificity / false-positive rate
    \item anti-degenerate check = Bumpy recall + F0.5
\end{itemize}

It is important to notice that we should not optimise precision alone and then hide recall. A model that would never predict `Bumpy`can obtain excellent precision, while being useless in practice.

\fbox{%
  \parbox{1.0\textwidth}{%
    \textbf{\underline{Technology Statement}}: During the preparation of this work, we used OpenAI ChatGPT (GPT-5.6 Sol), Google Gemini (3.1 Pro) and Claude Sonnet 5 Medium in order to brainstorm assignment direction, support on coding and discuss specific technical concepts to support the team along the development. The following parts of the assignment were affected/generated by AI tool usage: [DATA PREPARATION / MODELING / EVALUATION]. After using this tool/service, [Su Acar, Vikram Chandrasekaran and Henrique da Silva Couto] evaluated the validity of the tool’s outputs, including the sources that generative AI tools have used, and edited the content as needed. As a consequence, [Su Acar, Vikram Chandrasekaran and Henrique da Silva Couto] take full responsibility for the content of their work.
  }%
}
